\section{The critical half power expansion}
\label{spb:ep:sec:critical}
In this section $p=1$, $q=2$, $c=\sqrt{e/2}$,
$t=n^{-1/2}$, and $\lambda=c/t$.
The normalization $M_1$ is \eqref{spb:ep:eq:criticalM}.
Since $\lambda^2/n=c^2$, the leading quadratic defect persists:
$\Delta_n(X)$ is approximately $-3(tX)^2/4$, whose limiting value is
$-3c^2/4$. Expanding $e^{\Delta_n(X)}$ as a small perturbation of one would
miss an order-one factor. We retain that factor and expand around the
Poisson center instead.

\begin{theorem}\label{spb:ep:thm:critical}
For every fixed integer $K\ge0$, there are explicit polynomials
$C_{1,j}\in\Q[c]$ such that
\begin{equation}\label{spb:ep:eq:criticalall}
 \frac{a_1(n)}{M_1(n)}
     =\sum_{j=0}^K C_{1,j}(c)n^{-j/2}
                    +O_K(n^{-(K+1)/2}).
\end{equation}
Both parities have the same coefficients. There are no odd quarter powers
in the normalized expansion. The first four nonconstant coefficients are
\begin{align}
 C_{1,1}(c)&=\frac c4+\frac{11c^3}{8},\label{spb:ep:eq:critC1}\\
 C_{1,2}(c)&=\frac{31c^2}{32}-\frac{319c^4}{96}
                                      +\frac{121c^6}{128},\label{spb:ep:eq:critC2}\\
 C_{1,3}(c)&=-\frac{5c}{96}-\frac{739c^3}{128}
              +\frac{46517c^5}{3840}-\frac{7381c^7}{1536}
              +\frac{1331c^9}{3072},\label{spb:ep:eq:critC3}\\
 C_{1,4}(c)&=-\frac{227c^2}{128}+\frac{61923c^4}{2048}
              -\frac{771061c^6}{15360}+\frac{4140367c^8}{184320}
              \nonumber\\[-2pt]
          &\hspace{18mm}-\frac{41261c^{10}}{12288}
                         +\frac{14641c^{12}}{98304}.\label{spb:ep:eq:critC4}
\end{align}
\end{theorem}

\subsection{A separate global bound}
The exact Poisson identity is still available, now in the form
\begin{equation}\label{spb:ep:eq:criticalexact}
 \frac{a_1(n)}{M_1(n)}
  =2e^{3c^2/4}\E\left[e^{\Delta_n(X)}\1_{\{X\le n\}}
                               \1_{\{X\equiv n\ (2)\}}\right].
\end{equation}
The nonpositivity statement of Lemma~\ref{spb:ep:lem:sign} is false at $p=1$;
for example, $\Delta_n(1)>0$.
The required replacement is a uniform upper bound.

\begin{lemma}\label{spb:ep:lem:criticalbound}
For every integer $0\le r\le n$,
\begin{equation}\label{spb:ep:eq:criticalbound}
 \Delta_n(r)\le-\frac{r^2}{4n}+\frac r{2n}
              =\frac{1-(r-1)^2}{4n}\le\frac1{4n}.
\end{equation}
\end{lemma}
\begin{proof}
The first two terms of \eqref{spb:ep:eq:Delta} now form
$F(r)=((n+r)/2)\log(1+r/n)-r/2$.
One has $F(0)=F'(0)=0$ and
$F''(r)=1/(2n(1+r/n))\le1/(2n)$, so $F(r)\le r^2/(4n)$.
Since $\log(1-x)\le-x$ for $0\le x<1$,
\[
 \sum_{h=0}^{r-1}\log\left(1-\frac{2h}{n+r}\right)
       \le-\frac{r(r-1)}{n+r}\le-\frac{r(r-1)}{2n}.
\]
The final inequality uses $r(r-1)\ge0$, which holds at every allowed
integer. Adding the bounds proves \eqref{spb:ep:eq:criticalbound}.
\end{proof}

Let
\begin{equation}\label{spb:ep:eq:criticalbulk}
 \mathcal C_n=\{|X-\lambda|\le\lambda^{3/4}\}.
\end{equation}
The usual exponential Markov argument for the Poisson generating function
implies, for a positive absolute constant $b_0$ and large $\lambda$,
\begin{equation}\label{spb:ep:eq:criticaltail}
 \Pp(\mathcal C_n^c)
 \le2\exp\left(-\frac{\lambda^{3/2}}
                         {2(\lambda+\lambda^{3/4})}\right)
 =O(e^{-b_0\sqrt\lambda}).
\end{equation}
For completeness, the upper-tail Chernoff rate is
$\lambda((1+s)\log(1+s)-s)$ at displacement $s\lambda$;
integrating $\log(1+s)\ge s/(1+s)$ bounds it below by
$\lambda s^2/(2(1+s))$. The lower-tail rate is at least
$\lambda s^2/2$. Substitute $s=\lambda^{-1/4}$.

Lemma~\ref{spb:ep:lem:criticalbound} bounds the exact normalized contribution
on $\mathcal C_n^c$ by a fixed multiple of \eqref{spb:ep:eq:criticaltail}.
Every fixed polynomially weighted tail is also negligible by
Cauchy--Schwarz and Poisson moment bounds. Since $\lambda\asymp\sqrt n$,
these errors are smaller than every power of $n^{-1}$.
For large $n$, the central event lies inside $0\le X\le2\lambda<n$.

\subsection{Defect polynomials at criticality}
Specializing \eqref{spb:ep:eq:Dgeneral} gives
\begin{equation}\label{spb:ep:eq:Dcrit}
 D_j(r)=(-1)^{j+1}\left\{
  \frac{r^{j+1}}{2j(j+1)}
  +\frac1j\sum_{h=0}^{r-1}\bigl[(r-2h)^j-r^j\bigr]\right\}.
\end{equation}
The first values are
\begin{align*}
 D_1&=-3r^2/4+r,& D_2&=r^3/4-r/3,\\
 D_3&=-7r^4/24+r^3/3,&
 D_4&=7r^5/40-r^3/3+2r/15,\\
 D_5&=-11r^6/60+r^5/5.
\end{align*}
We require a local logarithmic estimate on a larger analytic scale than
in Lemma~\ref{spb:ep:lem:cauchy}.
For every fixed $L\ge1$, uniformly for integers $0\le r\le2\lambda$,
\begin{equation}\label{spb:ep:eq:criticalDrem}
 \Delta_n(r)=\sum_{j=1}^L D_j(r)n^{-j}
          +O_L\left(\frac{r^{L+2}}{n^{L+1}}\right)
       =\sum_{j=1}^L t^{2j}D_j(r)+O_L(t^L).
\end{equation}
To verify this, rewrite the defect exactly as
\[
 \Delta_n(r)=\frac{n-r}{2}\log(1+r/n)-\frac r2
            +\sum_{h=0}^{r-1}\log(1+(r-2h)/n).
\]
Expand the first logarithm through degree $L+1$ and each logarithm in
the sum through degree $L$. Their arguments are within $O(t)$ of one.
The usual logarithmic remainders, together with the $n,r$ prefactors
and the $r$ terms in the sum, give $O(r^{L+2}/n^{L+1})$.
The statement at $r=0$ is exact. Collecting coefficients gives
\eqref{spb:ep:eq:Dcrit}.

\subsection{Centered moments and a finite weighted generator}
Put
\begin{equation}\label{spb:ep:eq:Ucenter}
 U=t(X-\lambda)=tX-c.
\end{equation}
Let $B_{d,k}^{\ge2}$ count set partitions of $d$ labeled elements into
$k$ blocks, with no singleton blocks. They satisfy
\begin{equation}\label{spb:ep:eq:assocStirling}
 B_{0,0}^{\ge2}=1,\qquad
 B_{d,k}^{\ge2}=kB_{d-1,k}^{\ge2}+(d-1)B_{d-2,k-1}^{\ge2},
\end{equation}
with out-of-range indices zero. The recurrence separates the block
containing the last label according as it has more than two elements or
has two elements.
The centered Poisson generating function is
\[
 \E e^{z(X-\lambda)}=\exp(\lambda(e^z-1-z)).
\]
Coefficient comparison gives the exact polynomial identity
\begin{equation}\label{spb:ep:eq:centralmoment}
 \E U^d=\sum_{k=0}^{\lfloor d/2\rfloor}
           B_{d,k}^{\ge2}c^k t^{d-k}.
\end{equation}
Thus only integer powers of $t$ occur, and their least possible exponent
is $\lceil d/2\rceil$. In particular,
$\E|U|^{2h}=O_h(t^h)$ for each fixed $h$.

For the chosen expansion order $K$, define
\begin{equation}\label{spb:ep:eq:VK}
 V_K(t,u)=\frac{3c^2}{4}
            +\sum_{j=1}^{K+1}t^{2j}D_j((c+u)/t).
\end{equation}
It is a polynomial with no negative powers of $t$ because
$\deg D_j\le j+1$. Its constant term vanishes and
\begin{equation}\label{spb:ep:eq:VK0}
 V_K(0,u)=-\frac{3cu}{2}-\frac{3u^2}{4}.
\end{equation}
Assign weight two to $t$ and weight one to $u$.
Let $P_K(t,u)$ be $e^{V_K(t,u)}$ truncated to weighted degree at most
$2K+1$. Every monomial of $V_K$ has positive weight, so it suffices to
truncate the finite sum
\begin{equation}\label{spb:ep:eq:PK}
 \sum_{m=0}^{2K+1}\frac{V_K(t,u)^m}{m!}.
\end{equation}
Replace every $u^d$ in $P_K$ by the right side of
\eqref{spb:ep:eq:centralmoment}, and retain powers through $t^K$.
This finite procedure defines $C_{1,0},\ldots,C_{1,K}$.
It is different from the small-deformation generator of the previous section.

\subsection{Remainder and parity proof}
\begin{proof}[Proof of Theorem~\ref{spb:ep:thm:critical}]
On $\mathcal C_n$, one has
$|U|\le t\lambda^{3/4}=O(t^{1/4})\to0$.
Take $L=K+1$ in \eqref{spb:ep:eq:criticalDrem}. The exponent error is
$O(t^{K+1})$, and $V_K(t,U)$ is uniformly bounded on this event.
Therefore
\begin{equation}\label{spb:ep:eq:criticalexprem}
 e^{\Delta_n(X)+3c^2/4}
       =e^{V_K(t,U)}+O_K(t^{K+1})
       \quad\hbox{on }\mathcal C_n.
\end{equation}

Set $s=\sqrt t$ solely to apply ordinary multivariate Taylor's theorem to
$e^{V_K(s^2,u)}$ at $(0,0)$. Its total-degree $2K+1$ Taylor polynomial is
$P_K(s^2,u)$. The remainder is
$O_K((|s|+|u|)^{2K+2})$ in a fixed neighborhood.
Averaging this bound, using \eqref{spb:ep:eq:centralmoment} and
$(a+b)^m\le2^{m-1}(a^m+b^m)$, gives $O_K(t^{K+1})$.
The parity-restricted expectation costs at most a factor two.
The extra odd weighted degree is important: a Taylor argument stopped at
$2K$ would give only an absolute $O(t^{K+1/2})$ bound by this method.

The polynomial tail estimate and \eqref{spb:ep:eq:criticaltail} allow us to replace
the central expectation of $P_K$ by its unrestricted Poisson expectation.
For every fixed polynomial $Q(X)$ the exact parity filter says
\[
 2\E[Q(X)\1_{\{X\equiv n\ (2)\}}]
    =\E Q(X)+(-1)^n\E[(-1)^XQ(X)].
\]
For $Q(X)=X^j$, the last expectation is
$e^{-2\lambda}T_j(-\lambda)$.
A monomial $t^iU^d=t^i(tX-c)^d$ is a finite combination of
$t^{i+j}X^j$. Since $t^jT_j(-c/t)$ remains bounded, the parity contribution
from $P_K$ is $O_K(e^{-2\lambda})$.
Finally \eqref{spb:ep:eq:centralmoment} produces only integer powers of $t$;
those beyond $K$ are $O_K(t^{K+1})$.
Combining these statements with \eqref{spb:ep:eq:criticalexact} proves the
expansion and the claimed absence of odd quarter powers.
The displayed coefficients result from the finite generator.
\end{proof}

\subsection{The first coefficient without symbolic expansion}
There is a short independent way to see the first correction.
Put $Z=tX=c+U$ and $f(z)=e^{-3z^2/4}$.
On the bulk, the first two defect polynomials give
\[
 e^{\Delta_n(X)}=f(Z)\{1+t(Z+Z^3/4)+O(t^2)\}.
\]
Since $\E U=0$ and $\E U^2=ct$, the centered Taylor calculation gives
\[
 \E f(Z)=f(c)+\frac{ct}{2}f''(c)+O(t^2).
\]
Here a fourth-order Taylor remainder and the exact third central moment
justify the $O(t^2)$ error; an absolute third-moment bound alone would be
insufficient. The same argument applied to $f(Z)(Z+Z^3/4)$ gives its
expectation equal to its value at $c$ plus $O(t)$.
The tails have already been controlled globally.
Dividing by $f(c)$, the coefficient is
\[
 c+\frac{c^3}{4}+\frac c2\left(-\frac32+\frac{9c^2}{4}\right)
       =\frac c4+\frac{11c^3}{8}.
\]
This identifies the missing cancellation without a symbolic black box.
It also explains why adding the observed extra $1/(8c)$ changes the true
refinement while leaving the leading equivalent untouched.
